% !TEX program = xelatex
\documentclass[14pt]{extarticle}

\usepackage{amsmath}
\usepackage[margin=2.5cm]{geometry}
\usepackage{graphicx}
\usepackage{booktabs}
\usepackage[labelfont={bf},font={normalsize}]{caption}
\usepackage{xepersian}
\settextfont{XB Niloofar}

\title{برازش پارامترهای کیهان‌شناسی از طیف توان قطعات آسمان (فقط \lr{TT})}
\author{}
\date{\today}

\begin{document}
\maketitle

\section*{خلاصه}
با استفاده از نقشه‌های دمایی \lr{SMICA} پلانک (\lr{hm1 $\times$ hm2})،
طیف متقاطع \lr{TT} و کوواریانس گاوسی آن با \lr{NaMaster} روی ماسک شدت مشترک
برآورد شد، و پارامترهای \lr{$\Lambda$CDM} به‌همراه دامنه چشمه‌های نقطه‌ای
$A_{ps}^{TT}$ با \lr{iminuit} برازش شدند، در حالی که $A_s$ و $n_s$ در
مقادیر مرجع پلانک ۲۰۱۸ ثابت نگه داشته شدند. همین پردازش 
(
تنها-\lr{TT}
)
سپس به‌طور مستقل روی هر یک از ۱۲ قطعه مجزای \lr{HEALPix} با
$N_{side}=1$ (شکل~\ref{fig:patches}) تکرار شد. جدول~\ref{tab:params} پارامترهای بهترین‌برازش و
عدم‌قطعیت‌های به‌دست‌آمده برای هر قطعه را فهرست می‌کند، و شکل‌های
\ref{fig:bandpowers} و \ref{fig:paramsplot} به ترتیب طیف‌توان‌های هر قطعه و
پارامترهای بازیابی‌شده در قطعات مختلف را نشان می‌دهند.

\begin{figure}[htbp]
    \centering
    \includegraphics[width=1\textwidth]{../output/patches_map.png}
    \caption{۱۲ محدوده‌های مجزای $N_{side}=1$ که با ماسک پلانک
    (اپودایز‌شده) هم‌پوشانی داده شده‌اند.}
    \label{fig:patches}
\end{figure}

\begin{table}[htbp]
    \centering
    \small
    \caption{پارامترهای بهترین‌برازش هر قطعه (فقط \lr{TT}؛ $A_s$ و $n_s$ در
    مقادیر مرجع پلانک ۲۰۱۸، به‌ترتیب $2.1\times10^{-9}$ و $0.9649$، ثابت
    نگه داشته شده‌اند).}
    \label{tab:params}
    \begin{tabular}{cccccccc}
        \toprule
        قطعه & $f_{sky}$ & $\chi^2$ & همگرایی & $H_0$ & $\Omega_b h^2$ & $\Omega_c h^2$ & $A_{ps}^{TT}$ \\
        \midrule
        $0$  & $0.072$ & $7.20$  & بله & $66.32 \pm 2.42$ & $0.02104 \pm 0.00074$ & $0.1191 \pm 0.0044$ & $74.0 \pm 19.1$ \\
        $1$  & $0.073$ & $7.20$  & بله & $65.74 \pm 0.06$ & $0.02210 \pm 0.00002$ & $0.1234 \pm 0.0001$ & $55.1 \pm 14.4$ \\
        $2$  & $0.078$ & $6.85$  & بله & $68.33 \pm 1.90$ & $0.02155 \pm 0.00062$ & $0.1148 \pm 0.0034$ & $74.1 \pm 17.8$ \\
        $3$  & $0.075$ & $7.51$  & بله & $66.78 \pm 2.29$ & $0.02258 \pm 0.00077$ & $0.1208 \pm 0.0041$ & $40.0 \pm 20.0$ \\
        $4$  & $0.023$ & $6.55$  & بله & $72.00 \pm 3.53$ & $0.02341 \pm 0.00122$ & $0.1097 \pm 0.0056$ & $59.8 \pm 34.4$ \\
        $5$  & $0.040$ & $3.42$  & بله & $67.16 \pm 2.51$ & $0.02131 \pm 0.00093$ & $0.1189 \pm 0.0048$ & $80.2 \pm 21.3$ \\
        $6$  & $0.034$ & $6.62$  & بله & $71.10 \pm 3.19$ & $0.02261 \pm 0.00108$ & $0.1129 \pm 0.0056$ & $79.7 \pm 29.5$ \\
        $7$  & $0.049$ & $6.56$  & بله & $69.68 \pm 2.31$ & $0.02215 \pm 0.00085$ & $0.1109 \pm 0.0041$ & $63.8 \pm 19.1$ \\
        $8$  & $0.076$ & $8.28$  & بله & $66.11 \pm 1.72$ & $0.02199 \pm 0.00067$ & $0.1230 \pm 0.0033$ & $62.6 \pm 17.8$ \\
        $9$  & $0.076$ & $11.39$ & خیر & $71.66 \pm 8.21$ & $0.02284 \pm 0.00091$ & $0.1109 \pm 0.0166$ & $58.5 \pm 22.0$ \\
        $10$ & $0.073$ & $3.99$  & بله & $66.59 \pm 2.72$ & $0.02162 \pm 0.00075$ & $0.1198 \pm 0.0051$ & $62.5 \pm 17.8$ \\
        $11$ & $0.070$ & $12.44$ & بله & $67.14 \pm 2.14$ & $0.02205 \pm 0.00069$ & $0.1199 \pm 0.0042$ & $58.6 \pm 16.4$ \\
        \bottomrule
    \end{tabular}
\end{table}

\begin{figure}[htbp]
    \centering
    \includegraphics[width=\textwidth]{../output/patched_bandpowers_TT_fixed(As,ns).png}
    \caption{طیف‌توان‌های \lr{TT} هر قطعه، $\ell(\ell+1)C_\ell^{TT}/2\pi$.}
    \label{fig:bandpowers}
\end{figure}

\begin{figure}[htbp]
    \centering
    \includegraphics[width=0.85\textwidth]{../output/patched_params_TT_fixed(As,ns).png}
    \caption{پارامترها با بهترین‌ برازش به‌طور مستقل برای هر قطعه،
    در مقایسه با کیهان‌شناسی مرجع پلانک ۲۰۱۸.}
    \label{fig:paramsplot}
\end{figure}

\end{document}
